% (c) Nikita Lisitsa, lisyarus@gmail.com, 2026

\documentclass[10pt]{beamer}

\usepackage[T2A]{fontenc}
\usepackage[russian]{babel}
\usepackage{minted}

\usepackage[most]{tcolorbox}
\tcbuselibrary{minted}

\usepackage{graphicx}
\graphicspath{ {./images/} }

\usepackage{adjustbox}

\usepackage{color}
\usepackage{soul}
\usepackage{multicol}
\usepackage{hyperref}
\usepackage{tikz}
\usetikzlibrary{arrows.meta,positioning,calc}

\usetheme{metropolis}

\definecolor{red}{rgb}{1,0,0}
\definecolor{codebg}{RGB}{29,35,49}
\setminted{fontsize=\footnotesize}

\makeatletter
\newcommand{\slideimage}[1]{
  \begin{figure}
    \begin{adjustbox}{width=\textwidth, totalheight=\textheight-2\baselineskip-2\baselineskip,keepaspectratio}
      \includegraphics{#1}
    \end{adjustbox}
  \end{figure}
}
\makeatother

\title{Язык программирования C++}
\subtitle{Лекция 21: type traits, SFINAE, enable\_if, concepts}
\date{2026}
\author{Никита Лисица}

\setbeamertemplate{footline}[frame number]

\begin{document}

\frame{\titlepage}

\begin{frame}[fragile]
\frametitle{Напоминание: iterator\_traits}
\begin{itemize}
\usemintedstyle{tango}
\item Стандартный шаблон \mintinline{cpp}|std::iterator_traits<It>| предоставляет информацию об итераторах: категория (random access / bidirectional / etc), тип значения (\mintinline{cpp}|value_type|), и т.п.
\pause
\item По умолчанию достаёт информацию из самого типа итератора \mintinline{cpp}|It|, но содержит специализации для указателей
\pause
\item Такие шаблоны, содержащие некоторую информацию о типах, называются \textit{type traits}
\pause
\item В заголовочном файле \mintinline{cpp}|<type_traits>| очень много полезных type traits
\end{itemize}
\end{frame}

\begin{frame}[fragile]
\frametitle{Type traits}
\begin{itemize}
\usemintedstyle{lightbulb}
\begin{tcblisting}{listing engine=minted, minted language=cpp, minted options={fontsize=\scriptsize}, listing only, colback=codebg, colframe=codebg, rounded corners}
template <typename T>
struct is_void {
  static const bool value = false;
};

template <>
struct is_void<void> {
  static const bool value = true;
};
\end{tcblisting}
\pause
\usemintedstyle{lightbulb}
\begin{tcblisting}{listing engine=minted, minted language=cpp, minted options={fontsize=\scriptsize}, listing only, colback=codebg, colframe=codebg, rounded corners}
template <typename T>
struct is_pointer {
  static const bool value = false;
};

template <typename T>
struct is_pointer<T*> {
  static const bool value = true;
};
\end{tcblisting}
\end{itemize}
\end{frame}

\begin{frame}[fragile]
\frametitle{Type traits}
\begin{itemize}
\usemintedstyle{tango}
\item Ещё есть \mintinline{cpp}|std::is_integral|, \mintinline{cpp}|std::is_floating_point|, \mintinline{cpp}|std::is_bool|, \mintinline{cpp}|std::is_class|, \mintinline{cpp}|std::is_function|, \mintinline{cpp}|std::is_enum|, и т.п.
\pause
\item Некоторые из них можно реализовать на чистом C++, но многие (напр. \mintinline{cpp}|std::is_enum|) реализуются через компилятор
\pause
\item Для всех таких traits есть вспомогательные шаблонные константы с суффиксом \texttt{\_v}, чтобы не писать \texttt{::value}
\usemintedstyle{lightbulb}
\begin{tcblisting}{listing engine=minted, minted language=cpp, minted options={fontsize=\scriptsize}, listing only, colback=codebg, colframe=codebg, rounded corners}
template <typename T>
static constexpr bool is_integral_v = is_integral<T>::value;

// То же самое, что и
// is_integral<int>::value
std::is_integral_v<int>
\end{tcblisting}
\end{itemize}
\end{frame}

\begin{frame}[fragile]
\frametitle{Type traits}
\begin{itemize}
\usemintedstyle{tango}
\item Есть вспомогательные функции на типах \mintinline{cpp}|std::add/remove_const|, \mintinline{cpp}|std::add/remove_pointer|, \mintinline{cpp}|std::make_signed/unsigned|, и т.п.
\pause
\item Их результат -- не значение, а тип: \mintinline{cpp}|std::make_signed<int>::type|
\pause
\item Для них есть вспомогательные шаблонные \mintinline{cpp}|using|'и с суффиксом \texttt{\_t}, чтобы не писать \texttt{::type}
\usemintedstyle{lightbulb}
\begin{tcblisting}{listing engine=minted, minted language=cpp, minted options={fontsize=\scriptsize}, listing only, colback=codebg, colframe=codebg, rounded corners}
template <typename T>
using make_unsigned_t = typename make_unsigned<T>::type;

// То же самое, что и
// make_unsigned<int>::type
std::make_unsigned_t<int>
\end{tcblisting}
\end{itemize}
\end{frame}

\begin{frame}[fragile]
\frametitle{Type traits: использование}
\begin{itemize}
\usemintedstyle{tango}
\item Можно проверять корректность типов аргументов:
\usemintedstyle{lightbulb}
\begin{tcblisting}{listing engine=minted, minted language=cpp, minted options={fontsize=\scriptsize}, listing only, colback=codebg, colframe=codebg, rounded corners}
template <typename T>
T pow(T x, T n) {
  static_assert(std::is_integral_v<T>);
  // ...
}
\end{tcblisting}
\pause
\item Можно специализировать структуры данных и алгоритмы:
\usemintedstyle{lightbulb}
\begin{tcblisting}{listing engine=minted, minted language=cpp, minted options={fontsize=\scriptsize}, listing only, colback=codebg, colframe=codebg, rounded corners}
template <typename T, bool IsUnsigned = std::is_unsigned_v<T>>
class PowImpl { ... };

template <typename T>
class PowImpl<T, false> { ... };
\end{tcblisting}
\pause
\item \textbf{NB}: и то, и другое лучше делать с помощью \textit{концептов}
\end{itemize}
\end{frame}

\begin{frame}[fragile]
\frametitle{SFINAE}
\begin{itemize}
\usemintedstyle{tango}
\item Допустим, у нас есть шаблонная функция с набором перегрузок:
\usemintedstyle{lightbulb}
\begin{tcblisting}{listing engine=minted, minted language=cpp, minted options={fontsize=\scriptsize}, listing only, colback=codebg, colframe=codebg, rounded corners}
template <typename T>
T sum(T x, T y) {
  return x + y;
}

template <typename T>
auto sum(typename T::iterator begin, typename T::iterator end) {
  decltype(*begin) result = {};
  for (auto it = begin; it != end; ++it)
    result += *it;
  return result;
}
\end{tcblisting}
\pause
\usemintedstyle{tango}
\item Что произойдёт, если вызвать \mintinline{cpp}|sum<int>(4, 5)|?
\end{itemize}
\end{frame}

\begin{frame}[fragile]
\frametitle{SFINAE}
\begin{itemize}
\usemintedstyle{tango}
\item Что произойдёт, если вызвать \mintinline{cpp}|sum<int>(4, 5)|? \pause У типа \mintinline{cpp}|int| нет вложенного типа \mintinline{cpp}|iterator|, поэтому должна произойти ошибка компиляции
\pause
\item Стандарт C++ говорит в таких случаях игнорировать те перегрузки, что вызывают \textit{ошибку (failure)} в процессе \textit{подстановки (substitution)} шаблонного параметра
\pause
\item \textbf{S}ubstitution \textbf{F}ailure \textbf{I}s \textbf{N}ot \textbf{A}n \textbf{E}rror, или \textbf{SFINAE}
\pause
\item \(\Longrightarrow\) Выберется другая перегрузка \mintinline{cpp}|sum(T, T)|
\end{itemize}
\end{frame}

\begin{frame}[fragile]
\frametitle{SFINAE: применение}
\begin{itemize}
\usemintedstyle{tango}
\item Можно узнать, есть ли у типа вложенный тип \mintinline{cpp}|iterator|:
\usemintedstyle{lightbulb}
\begin{tcblisting}{listing engine=minted, minted language=cpp, minted options={fontsize=\scriptsize}, listing only, colback=codebg, colframe=codebg, rounded corners}
template <typename T>
struct has_iterator {
  template <typename U>
  static char test(typename U::iterator*);

  template <typename U>
  static long test(...);

  static constexpr bool value =
    (sizeof(test<T>(nullptr)) == sizeof(char));
};

has_iterator<std::vector<int>>::value; // == true
\end{tcblisting}
\end{itemize}
\end{frame}

\begin{frame}[fragile]
\frametitle{SFINAE: применение}
\begin{itemize}
\usemintedstyle{tango}
\item Можно проверить, если ли перегрузка некоторой функции от типа:
\usemintedstyle{lightbulb}
\begin{tcblisting}{listing engine=minted, minted language=cpp, minted options={fontsize=\scriptsize}, listing only, colback=codebg, colframe=codebg, rounded corners}
template <typename T>
struct is_printable {
  template <typename U>
  static char test(decltype(print(U))*);

  template <typename U>
  static long test(...);

  static constexpr bool value =
    (sizeof(test<T>(nullptr)) == sizeof(char));
};
\end{tcblisting}
\end{itemize}
\end{frame}

\begin{frame}[fragile]
\frametitle{SFINAE: применение}
\begin{itemize}
\usemintedstyle{tango}
\item Можно использовать напрямую, без классов-обёрток:
\usemintedstyle{lightbulb}
\begin{tcblisting}{listing engine=minted, minted language=cpp, minted options={fontsize=\scriptsize}, listing only, colback=codebg, colframe=codebg, rounded corners}
template <typename T>
void show(const T& x, typename T::iterator* = nullptr) {
  std::cout << x;
}
\end{tcblisting}
\end{itemize}
\end{frame}

\begin{frame}[fragile]
\frametitle{enable\_if}
\begin{itemize}
\usemintedstyle{tango}
\item \mintinline{cpp}|std::enable_if| -- шаблон, помогающий за счёт SFINAE специализировать функции по какому-либо условию
\pause
\usemintedstyle{lightbulb}
\begin{tcblisting}{listing engine=minted, minted language=cpp, minted options={fontsize=\scriptsize}, listing only, colback=codebg, colframe=codebg, rounded corners}
template <bool Condition, typename T>
struct enable_if {
  using type = T;
};

template <typename T>
struct enable_if<false, T> {};

template <typename T>
using enable_if_t = typename enable_if<T>::type;
\end{tcblisting}
\end{itemize}
\end{frame}

\begin{frame}[fragile]
\frametitle{enable\_if}
\begin{itemize}
\usemintedstyle{lightbulb}
\begin{tcblisting}{listing engine=minted, minted language=cpp, minted options={fontsize=\scriptsize}, listing only, colback=codebg, colframe=codebg, rounded corners}
template <typename T>
typename std::enable_if_t<!has_iterator<T>, void>
print(const T& x) {
  std::cout << x;
}

template <typename T>
typename std::enable_if_t<has_iterator<T>, void>
void print(const T& x) {
  for (const auto& v : x)
    std::cout << v;
}
\end{tcblisting}
\end{itemize}
\end{frame}

\begin{frame}[fragile]
\frametitle{Концепты (concepts)}
\begin{itemize}
\usemintedstyle{tango}
\item Появились в C++20
\pause
\item Изначально задумывались как сильное расширение системы типов C++ в сторону Haskell typeclasses и higher-order types
\pause
\item В итоге свелись к аналогу \texttt{SFINAE/enable\_if}, но с нормальным синтаксисом и с поддержкой со стороны компилятора
\pause
\item \textbf{NB}: В первой версии концептов можно было задавать \textit{аксиомы} на операции и типы (например, что сложение ассоциативно)
\end{itemize}
\end{frame}

\begin{frame}[fragile]
\frametitle{Концепты (concepts)}
\begin{itemize}
\usemintedstyle{tango}
\item Сам концепт -- по сути, булевая константа, параметризованная типом, известная на момент компиляции
\pause
\usemintedstyle{lightbulb}
\begin{tcblisting}{listing engine=minted, minted language=cpp, minted options={fontsize=\scriptsize}, listing only, colback=codebg, colframe=codebg, rounded corners}
template <typename T>
concept Integral = std::is_integral_v<T>;

template <typename T>
concept Unsigned = Integral<T> && std::is_unsigned_v<T>;
\end{tcblisting}
\end{itemize}
\end{frame}

\begin{frame}[fragile]
\frametitle{Requires expression}
\begin{itemize}
\usemintedstyle{tango}
\item Для упрощения задания концептов есть особый синтаксис: \textit{requires expression}
\pause
\usemintedstyle{lightbulb}
\begin{tcblisting}{listing engine=minted, minted language=cpp, minted options={fontsize=\scriptsize}, listing only, colback=codebg, colframe=codebg, rounded corners}
template <typename T>
concept Comparable = requires (T a, T b) {
  a < b; // Проверяет, что есть подходящий оператор <
};
\end{tcblisting}
\pause
\item Можно задавать несколько выражений, и проверять их типы:
\usemintedstyle{lightbulb}
\begin{tcblisting}{listing engine=minted, minted language=cpp, minted options={fontsize=\scriptsize}, listing only, colback=codebg, colframe=codebg, rounded corners}
template <typename T>
concept Comparable = requires (T a, T b) {
  { a < b } -> std::same_as<bool>;
  // или
  { a < b } -> std::convertible_to<bool>;
};
\end{tcblisting}
\end{itemize}
\end{frame}

\begin{frame}[fragile]
\frametitle{Requires expression}
\begin{itemize}
\usemintedstyle{lightbulb}
\begin{tcblisting}{listing engine=minted, minted language=cpp, minted options={fontsize=\scriptsize}, listing only, colback=codebg, colframe=codebg, rounded corners}
template <typename It>
concept ForwardIterator = requires (It it) {
  typename It::value_type;
  { *it } -> std::same_as<typename It::value_type&>;
  { ++it } -> std::same_as<It&>;
};
\end{tcblisting}
\end{itemize}
\end{frame}

\begin{frame}[fragile]
\frametitle{Использование концептов, requires clause}
\begin{itemize}
\usemintedstyle{tango}
\item Можно ограничить параметры шаблонной функции только теми, что удовлетворяют опеределённому концепту:
\usemintedstyle{lightbulb}
\begin{tcblisting}{listing engine=minted, minted language=cpp, minted options={fontsize=\scriptsize}, listing only, colback=codebg, colframe=codebg, rounded corners}
template <typename T>
requires Comparable<T>
T min(T x, T y) {
  return (x < y) ? x : y;
}
\end{tcblisting}
\pause
\item Можно комбинировать несколько концептов:
\usemintedstyle{lightbulb}
\begin{tcblisting}{listing engine=minted, minted language=cpp, minted options={fontsize=\scriptsize}, listing only, colback=codebg, colframe=codebg, rounded corners}
template <typename T>
requires Comparable<T> && Copyable<T>
T min(T x, T y) {
  return (x < y) ? x : y;
}
\end{tcblisting}
\end{itemize}
\end{frame}

\begin{frame}[fragile]
\frametitle{Использование концептов, requires clause}
\begin{itemize}
\usemintedstyle{tango}
\item Есть упрощённый синтаксис для простого случая одного концепта:
\usemintedstyle{lightbulb}
\begin{tcblisting}{listing engine=minted, minted language=cpp, minted options={fontsize=\scriptsize}, listing only, colback=codebg, colframe=codebg, rounded corners}
template <Comparable T>
T min(T x, T y) {
  return (x < y) ? x : y;
}
\end{tcblisting}
\pause
\item ...но его можно комбинировать и с другими концептами:
\usemintedstyle{lightbulb}
\begin{tcblisting}{listing engine=minted, minted language=cpp, minted options={fontsize=\scriptsize}, listing only, colback=codebg, colframe=codebg, rounded corners}
// То же самое, что и Copyable + Comparable
template <Comparable T>
requires Copyable<T>
T min(T x, T y) {
  return (x < y) ? x : y;
}
\end{tcblisting}
\end{itemize}
\end{frame}

\begin{frame}[fragile]
\frametitle{Использование концептов, requires clause}
\begin{itemize}
\usemintedstyle{tango}
\item То же самое можно использовать для шаблонов классов:
\usemintedstyle{lightbulb}
\begin{tcblisting}{listing engine=minted, minted language=cpp, minted options={fontsize=\scriptsize}, listing only, colback=codebg, colframe=codebg, rounded corners}
template <typename T>
requires Comparable<T>
struct sorted_array {
  // ...
};
\end{tcblisting}
\pause
\begin{tcblisting}{listing engine=minted, minted language=cpp, minted options={fontsize=\scriptsize}, listing only, colback=codebg, colframe=codebg, rounded corners}
template <Comparable T>
struct sorted_array {
  // ...
};
\end{tcblisting}
\end{itemize}
\end{frame}

\begin{frame}[fragile]
\frametitle{Requires requires}
\begin{itemize}
\usemintedstyle{tango}
\item Конструкция \mintinline{cpp}|requires|, позволяющая проверить набор условий на тип, называется \textit{requires expression}:
\usemintedstyle{lightbulb}
\begin{tcblisting}{listing engine=minted, minted language=cpp, minted options={fontsize=\scriptsize}, listing only, colback=codebg, colframe=codebg, rounded corners}
requires (T a) { a < b; }
\end{tcblisting}
\pause
\item Конструкция \mintinline{cpp}|requires|, проверяющая шаблонный параметр, называется \textit{requires clause}:
\usemintedstyle{lightbulb}
\begin{tcblisting}{listing engine=minted, minted language=cpp, minted options={fontsize=\scriptsize}, listing only, colback=codebg, colframe=codebg, rounded corners}
...
requires Comparable<T>
...
\end{tcblisting}
\pause
\item Что, если мы хотим определить безымянный концепт и сразу его использовать?
\end{itemize}
\end{frame}

\begin{frame}[fragile]
\frametitle{Requires requires}
\begin{itemize}
\usemintedstyle{tango}
\item Что, если мы хотим определить безымянный концепт и сразу его использовать? \pause \(\Longrightarrow\) \textit{requires requires}:
\usemintedstyle{lightbulb}
\begin{tcblisting}{listing engine=minted, minted language=cpp, minted options={fontsize=\scriptsize}, listing only, colback=codebg, colframe=codebg, rounded corners}
template <typename T>
requires requires (T x, T y) { x < y; }
T min(T x, T y) {
  return (x < y) ? x : y;
}
\end{tcblisting}
\pause
\usemintedstyle{tango}
\item Выражение \mintinline{cpp}|requires (T x, T y) { x < y; }| целиком -- compile-time булево значение (requires expression)
\pause
\item Выражение \mintinline{cpp}|requires Condition| ожидает compile-time булево значение в качестве параметра (requires clause)
\end{itemize}
\end{frame}

\begin{frame}[fragile]
\frametitle{Концепты (concepts)}
\begin{itemize}
\usemintedstyle{tango}
\item Чем концепты лучше, чем SFINAE + enable\_if?
\pause
\item Быстрее компиляция
\pause
\item Чище и понятнее код
\pause
\item Не полагаемся на хаки языка
\end{itemize}
\end{frame}

\begin{frame}[fragile]
\frametitle{Концепты (concepts)}
\begin{itemize}
\usemintedstyle{tango}
\item Зачем использовать концепты?
\pause
\item Специализация структур данных или алгоритмов по наборам типов
\pause
\item Улучшение ошибок компиляции
\end{itemize}
\end{frame}

\begin{frame}[fragile]
\frametitle{Ошибки компиляции в шаблонах}
\begin{itemize}
\usemintedstyle{tango}
\item Один из недостатков шаблонов как вида полиморфизма -- сложные ошибки компиляции
\pause
\item Пример: пытаемся отсортировать тип, у которого нет оператора сравнения
\usemintedstyle{lightbulb}
\begin{tcblisting}{listing engine=minted, minted language=cpp, minted options={fontsize=\scriptsize}, listing only, colback=codebg, colframe=codebg, rounded corners}
#include <algorithm>
#include <iterator>

struct point {
  int x;
  int y;
};

int main()
{
  point a[3] { {1,2}, {3,4}, {5,6} };
  std::sort(std::begin(a), std::end(a));
}
\end{tcblisting}
\pause
\usemintedstyle{tango}
\item Ожидание: `у класса \mintinline{cpp}|point| нет оператора \mintinline{cpp}|<|'
\end{itemize}
\end{frame}

\begin{frame}[fragile]
\frametitle{Ошибки компиляции в шаблонах}
\begin{itemize}
\usemintedstyle{tango}
\item Реальность:
\slideimage{compile_error.png}
\end{itemize}
\end{frame}

\begin{frame}[fragile]
\frametitle{Ошибки компиляции в шаблонах}
\begin{itemize}
\usemintedstyle{tango}
\item GCC 15: ошибка на 131 строку
\pause
\item \mintinline{cpp}|std::sort| вызывает \mintinline{cpp}|std::__sort|, который вызывает \mintinline{cpp}|std::__introsort_loop|, который вызывает \mintinline{cpp}|std::__partial_sort|, который вызывает \mintinline{cpp}|std::__heap_select|, который вызывает \mintinline{cpp}|std::__make_heap|, который вызывает \mintinline{cpp}|std::__adjust_heap|, который вызывает \mintinline{cpp}|std::__push_heap|, который вызывает сравнение объектов типа \mintinline{cpp}|point|
\pause
\item У компилятора нет информации о том, какой из приведённых вызовов на самом деле ошибочный, поэтому ему приходится выдавать весь стек вызова
\pause
\item Компилятор не знает, какой ожидался оператор, поэтому ему приходится вывести список всех опробованных операторов, включая абсолютно нерелевантные (например, для итераторов по \mintinline{cpp}|std::vector|, или для строк)
\end{itemize}
\end{frame}

\begin{frame}[fragile]
\frametitle{Ошибки компиляции в концептах}
\begin{itemize}
\usemintedstyle{tango}
\item Добавим в \mintinline{cpp}|std::sort| новый концепт \mintinline{cpp}|Sortable|:
\usemintedstyle{lightbulb}
\begin{tcblisting}{listing engine=minted, minted language=cpp, minted options={fontsize=\scriptsize}, listing only, colback=codebg, colframe=codebg, rounded corners}
template <typename It>
concept Sortable = requires (It x, It y) {
  { *x < *y } -> std::same_as<bool>;
};

template <Sortable It>
void sort(It begin, It end) {
  // ...
}
\end{tcblisting}
\end{itemize}
\end{frame}

\begin{frame}[fragile]
\frametitle{Ошибки компиляции в концептах}
\slideimage{concepts_error.png}
\end{frame}

\end{document}